\begin{claim}{B.52}{Open}{Exact remaining scope of the henselian-field question B8}
The classification demanded by B8 remains open for general
finite-transcendence-degree fields over $\mathbb Q$ with a
henselian mixed-characteristic valuation outside the following
proved positive and negative classes. All positive assertions
below concern full closed-point-preserving relative etale
reconstruction of finite-type schemes, with the structural
isomorphism retained and universal homeomorphisms inverted.
They include every finite extension of the stated field.

For every prime $p$, let
$K_p=\operatorname{Frac}(\mathbb Z_{(p)}^{sh})$ and $I=G_{K_p}$.
The fields
\[
 B_p=\overline{\mathbb Q}^{I''},\qquad
 K_{p,\mathrm{nil}}=
   \overline{\mathbb Q}^{\,\bigcap_{c\ge1}\gamma_{c+1}(I)}
\]
are positive by B.136 and B.138, respectively, including $p=2,3$.
Every fixed bounded-nilpotency field is also covered by B.123.

More generally, let $D\subset\overline{\mathbb Q}_p$ be
finitely generated over $\mathbb Q$, let $M$ be its finite
local completion, and let $\Omega$ be its algebraic closure
inside $\overline{\mathbb Q}_p$.
B.139--B.140 prove the corresponding finite-derived and
lower-central-union assertions in every finite transcendence
degree. B.142--B.143 prove positivity for
$\Omega^N$ for every nontrivial closed normal subgroup
$N\triangleleft G_M$, and for every algebraic subfield
over which such a positive field is algebraic.
A nontrivial normal core of the field's absolute Galois
subgroup in some open local Galois group is consequently
a sufficient criterion, with the relative local base
chosen as in B.139. No normal closure is substituted
for that core.

There is a further positive class with no finite-residue
restriction. Let $D/\mathbb Q$ be finitely generated with
a discrete mixed-characteristic valuation whose residue
field is finitely generated over $\mathbb F_p$.
Let $M$ be its completion, let $\Omega$ be the algebraic
closure of $D$ in $\overline M$, and let $I=I_M\subset G_D$.
B.146--B.147 prove positivity for
\[
 \Omega^{I^{(r)}}\quad(r\ge0),\qquad
 \Omega^{\,\bigcap_{c\ge1}\gamma_{c+1}(I)},
\]
and algebraic descent as above.
The Gauss valuations on $\mathbb Q(t_1,\ldots,t_s)$
give strictly additional families with residue
transcendence degree $s$ for every $s>0$.
The all-proper-normal-subgroup conclusion of B.143
is not asserted for these more general complete bases.

On the negative side, B.144 proves that positivity
requires a slim absolute Galois group: every open
subgroup must have trivial centre.
It gives explicit henselian fields of every finite
transcendence degree, and all their finite extensions,
whose absolute Galois groups are wild procyclic
$\mathbb Z_p$ groups. Their cores in every open
finite-local Galois group are trivial.
Reconstruction is also false for every field in the
stated finite-transcendence-degree setting admitting
a nontrivial henselian valuation of residue
characteristic zero, by B.49 and B.54--B.55.
This includes a mixed-characteristic valuation with
such a henselian coarsening.
An algebraically closed field is negative by B.53;
in particular the union of all finite derived fields
in the finite-local setting is distinguished from
the positive lower-central union.

Thus the remaining finite-local-subfield cases are
slim and have trivial normal core in every open
local group, apart from any further positive cases
already covered by the stated criteria.
Outside that setting, other algebraic Galois
subgroups, residue-field configurations and
nondiscrete valuations remain in the quantifier.
Neither unbounded wild ramification nor a
non-torsion Kummer kernel alone decides it.
\end{claim}
\begin{proof}[Scope verification and unresolved step]
B.136 proves the arbitrary-source norm-polynomial step
at every prime. B.137--B.138 verify the uniform two-grid
and local normalization inputs for the entire
lower-central union, including $2,3$.
Their final Carlson--Haine--Wolf reduction reconstructs
the full relative inverse-image functor.

B.139 proves the finite-extension and completion
comparison for the relative finitely generated base.
B.140 checks the finite-derived and nilpotent stable
Kummer kernels, two-grid rigidity, arbitrary-source
norm recovery and algebraic descent in that setting.
B.142 supplies the more general abelian normal
quotient criterion. B.143 proves that every
nontrivial normal subgroup of a finite-local
absolute Galois group is nonabelian; a least strict
derived step supplies the criterion. Algebraic
descent gives the asserted normal-core consequence.
All of these claims include finite extensions.

B.146 treats a complete discrete valuation field
with finitely generated, possibly imperfect residue
field. Its additive logarithmic map injects the
complete field into rational Kummer cohomology;
it is not asserted to be onto.
B.147 uses ratios in deep units to place the exact
scalar error in this subspace. It verifies every
remaining norm and graph input without assuming
local compactness or freeness of the general wild
inertia group. Its final algebraic descent and
two-henselian-valuations argument establish the
claimed additional scope.

For the negative procyclic fields, B.144 constructs
a nontrivial central absolute Galois element and
uses its inverse-image natural transformation to
produce an actual-point map on the affine line
which no polynomial realizes. Its finite-etale
slice argument shows that positivity would pass
to every finite extension, proving the slimness
necessity and the stated negative family.
The residue-zero conclusions use the arithmetic
covering-category comparison and valuation
approximation of B.49 and B.54--B.55; those
residue-characteristic hypotheses remain essential.

For a general remaining mixed-characteristic field,
the norm proof still requires a group controlling
normalized norm equivariance, arbitrarily large
cyclic quotients with the stated normal-subgroup
surjectivity, the scalar logarithmic comparison,
and a stable finite-grid rigidity test.
These inputs have now been verified for all
positive classes above, but not for every field
in the remaining quantifier.
The displayed negative constructions do not
decide those residual cases either.
The full classification therefore remains Open.
\end{proof}
